% Options for packages loaded elsewhere
\PassOptionsToPackage{unicode}{hyperref}
\PassOptionsToPackage{hyphens}{url}
\documentclass[
  12pt,
]{ctexart}
\usepackage{xcolor}
\usepackage{amsmath,amssymb}
\setcounter{secnumdepth}{-\maxdimen} % remove section numbering
\usepackage{iftex}
\ifPDFTeX
  \usepackage[T1]{fontenc}
  \usepackage[utf8]{inputenc}
  \usepackage{textcomp} % provide euro and other symbols
\else % if luatex or xetex
  \usepackage{unicode-math} % this also loads fontspec
  \defaultfontfeatures{Scale=MatchLowercase}
  \defaultfontfeatures[\rmfamily]{Ligatures=TeX,Scale=1}
\fi
\usepackage{lmodern}
\ifPDFTeX\else
  % xetex/luatex font selection
\fi
% Use upquote if available, for straight quotes in verbatim environments
\IfFileExists{upquote.sty}{\usepackage{upquote}}{}
\IfFileExists{microtype.sty}{% use microtype if available
  \usepackage[]{microtype}
  \UseMicrotypeSet[protrusion]{basicmath} % disable protrusion for tt fonts
}{}
\makeatletter
\@ifundefined{KOMAClassName}{% if non-KOMA class
  \IfFileExists{parskip.sty}{%
    \usepackage{parskip}
  }{% else
    \setlength{\parindent}{0pt}
    \setlength{\parskip}{6pt plus 2pt minus 1pt}}
}{% if KOMA class
  \KOMAoptions{parskip=half}}
\makeatother
\usepackage{longtable,booktabs,array}
\usepackage{calc} % for calculating minipage widths
% Correct order of tables after \paragraph or \subparagraph
\usepackage{etoolbox}
\makeatletter
\patchcmd\longtable{\par}{\if@noskipsec\mbox{}\fi\par}{}{}
\makeatother
% Allow footnotes in longtable head/foot
\IfFileExists{footnotehyper.sty}{\usepackage{footnotehyper}}{\usepackage{footnote}}
\makesavenoteenv{longtable}
\setlength{\emergencystretch}{3em} % prevent overfull lines
\providecommand{\tightlist}{%
  \setlength{\itemsep}{0pt}\setlength{\parskip}{0pt}}
\usepackage{bookmark}
\IfFileExists{xurl.sty}{\usepackage{xurl}}{} % add URL line breaks if available
\urlstyle{same}
\hypersetup{
  hidelinks,
  pdfcreator={LaTeX via pandoc}}

\author{SCX}
\date{2026}

\begin{document}

\section{Unified Notation and Theorem Dependency
Graph}\label{unified-notation-and-theorem-dependency-graph}

\begin{quote}
This document consolidates notation across all three SCX theorems and
maps the logical dependencies between them. It serves as a
quick-reference companion to the polished theorem statements.
\end{quote}

\begin{center}\rule{0.5\linewidth}{0.5pt}\end{center}

\subsection{1 Unified Notation Table}\label{unified-notation-table}

\subsubsection{1.1 Core Objects}\label{core-objects}

\begin{longtable}[]{@{}
  >{\raggedright\arraybackslash}p{(\linewidth - 10\tabcolsep) * \real{0.1739}}
  >{\raggedright\arraybackslash}p{(\linewidth - 10\tabcolsep) * \real{0.1957}}
  >{\raggedright\arraybackslash}p{(\linewidth - 10\tabcolsep) * \real{0.1522}}
  >{\raggedright\arraybackslash}p{(\linewidth - 10\tabcolsep) * \real{0.1522}}
  >{\raggedright\arraybackslash}p{(\linewidth - 10\tabcolsep) * \real{0.1522}}
  >{\raggedright\arraybackslash}p{(\linewidth - 10\tabcolsep) * \real{0.1739}}@{}}
\toprule\noalign{}
\begin{minipage}[b]{\linewidth}\raggedright
Symbol
\end{minipage} & \begin{minipage}[b]{\linewidth}\raggedright
Meaning
\end{minipage} & \begin{minipage}[b]{\linewidth}\raggedright
Thm 1
\end{minipage} & \begin{minipage}[b]{\linewidth}\raggedright
Thm 2
\end{minipage} & \begin{minipage}[b]{\linewidth}\raggedright
Thm 3
\end{minipage} & \begin{minipage}[b]{\linewidth}\raggedright
Domain
\end{minipage} \\
\midrule\noalign{}
\endhead
\bottomrule\noalign{}
\endlastfoot
\(\mathcal{X}\) & Input space & x & x & x & Measurable space \\
\(\mathcal{Y}\) & Label space, \(K = \|\mathcal{Y}\|\) & x & x & x &
Discrete (classif.) or \(\mathbb{R}^d\) (regr.) \\
\(X\) & Input random variable & x & x & x & \(\mathcal{X}\) \\
\(Y\) & Observed label & x & x & x & \(\mathcal{Y}\) \\
\(Y^*\) & True (latent) label & --- & --- & x & \(\mathcal{Y}\) \\
\(f^* : \mathcal{X} \to \mathcal{Y}\) & Ground-truth oracle & x & --- &
x & --- \\
\(\{f_m\}_{m=1}^M\) & Expert models & x & x & x &
\(\mathcal{X} \to \mathcal{Y}\) \\
\(\ell : \mathcal{Y} \times \mathcal{Y} \to [0, B]\) & Bounded loss
function & x & x & --- & \(B < \infty\) \\
\(\tau > 0\) & Expert error threshold & x & --- & --- &
\(\mathbb{R}^+\) \\
\end{longtable}

\subsubsection{1.2 Noise Model}\label{noise-model}

\begin{longtable}[]{@{}
  >{\raggedright\arraybackslash}p{(\linewidth - 10\tabcolsep) * \real{0.1600}}
  >{\raggedright\arraybackslash}p{(\linewidth - 10\tabcolsep) * \real{0.1800}}
  >{\raggedright\arraybackslash}p{(\linewidth - 10\tabcolsep) * \real{0.1400}}
  >{\raggedright\arraybackslash}p{(\linewidth - 10\tabcolsep) * \real{0.1400}}
  >{\raggedright\arraybackslash}p{(\linewidth - 10\tabcolsep) * \real{0.1400}}
  >{\raggedright\arraybackslash}p{(\linewidth - 10\tabcolsep) * \real{0.2400}}@{}}
\toprule\noalign{}
\begin{minipage}[b]{\linewidth}\raggedright
Symbol
\end{minipage} & \begin{minipage}[b]{\linewidth}\raggedright
Meaning
\end{minipage} & \begin{minipage}[b]{\linewidth}\raggedright
Thm 1
\end{minipage} & \begin{minipage}[b]{\linewidth}\raggedright
Thm 2
\end{minipage} & \begin{minipage}[b]{\linewidth}\raggedright
Thm 3
\end{minipage} & \begin{minipage}[b]{\linewidth}\raggedright
Definition
\end{minipage} \\
\midrule\noalign{}
\endhead
\bottomrule\noalign{}
\endlastfoot
\(\eta\) & Global noise rate \(\mathbb{P}(Y \neq Y^*)\) & x & x & x &
\(\eta \in (0, 1/2)\) \\
\(\eta(x)\) & Input-dependent noise rate & --- & --- & x &
\(\eta(x) = \mathbb{P}(Y \neq Y^* \mid X = x)\) \\
\(Z\) & Noise indicator (\(1\) = noise, \(0\) = clean) & --- & x & --- &
\(Z \in \{0, 1\}\) \\
\(e_m(x, y)\) & Expert error indicator & x & --- & --- &
\(e_m = \mathbf{1}\{\ell(f_m(x), y) > \tau\}\) \\
\(C(x)\) & Consistency score & x & --- & --- &
\(C(x) = \frac{1}{M}\sum_m e_m(x, y)\) \\
\end{longtable}

\subsubsection{1.3 State and Feature
Structure}\label{state-and-feature-structure}

\begin{longtable}[]{@{}
  >{\raggedright\arraybackslash}p{(\linewidth - 10\tabcolsep) * \real{0.1600}}
  >{\raggedright\arraybackslash}p{(\linewidth - 10\tabcolsep) * \real{0.1800}}
  >{\raggedright\arraybackslash}p{(\linewidth - 10\tabcolsep) * \real{0.1400}}
  >{\raggedright\arraybackslash}p{(\linewidth - 10\tabcolsep) * \real{0.1400}}
  >{\raggedright\arraybackslash}p{(\linewidth - 10\tabcolsep) * \real{0.1400}}
  >{\raggedright\arraybackslash}p{(\linewidth - 10\tabcolsep) * \real{0.2400}}@{}}
\toprule\noalign{}
\begin{minipage}[b]{\linewidth}\raggedright
Symbol
\end{minipage} & \begin{minipage}[b]{\linewidth}\raggedright
Meaning
\end{minipage} & \begin{minipage}[b]{\linewidth}\raggedright
Thm 1
\end{minipage} & \begin{minipage}[b]{\linewidth}\raggedright
Thm 2
\end{minipage} & \begin{minipage}[b]{\linewidth}\raggedright
Thm 3
\end{minipage} & \begin{minipage}[b]{\linewidth}\raggedright
Definition
\end{minipage} \\
\midrule\noalign{}
\endhead
\bottomrule\noalign{}
\endlastfoot
\(\mathcal{S}\) & State space & x & x & x & Finite set,
\(\|\mathcal{S}\| = K_S\) \\
\(s : \mathcal{X} \to \mathcal{S}\) & State assignment & x & x & x &
Partition of \(\mathcal{X}\) \\
\(\rho_s = \mathbb{P}(X \in s)\) & State probability & x & --- & x &
\(\sum_s \rho_s = 1\) \\
\(\mu_s\) & State-level clean error bound & x & --- & --- &
\(\mu_s = \sup_{x \in s} \mathbb{E}[C \mid \text{clean}, x]\) \\
\(C_{\text{bal}}\) & Error balance constant & x & --- & --- &
\(C_{\text{bal}} \geq 1\) \\
\(\Delta_s\) & Separation gap for state \(s\) & x & --- & --- &
\(\Delta_s = \min(\theta - \mu_s, 1 - C_{\text{bal}}\mu_s/(K-1) - \theta)\) \\
\(\phi : \mathcal{X} \to \Phi\) & Observed feature map & --- & x & --- &
\(\Phi \subseteq \mathbb{R}^{d_\phi}\) \\
\(\delta\) & Feature-state mutual information & --- & x & --- &
\(\delta = I(\phi(X); S)\) \\
\(\varepsilon_\phi\) & Normalized weakness & --- & x & --- &
\(\varepsilon_\phi = \delta / \log K_S \in [0, 1]\) \\
\end{longtable}

\subsubsection{1.4 Detection and
Evaluation}\label{detection-and-evaluation}

\begin{longtable}[]{@{}
  >{\raggedright\arraybackslash}p{(\linewidth - 10\tabcolsep) * \real{0.1600}}
  >{\raggedright\arraybackslash}p{(\linewidth - 10\tabcolsep) * \real{0.1800}}
  >{\raggedright\arraybackslash}p{(\linewidth - 10\tabcolsep) * \real{0.1400}}
  >{\raggedright\arraybackslash}p{(\linewidth - 10\tabcolsep) * \real{0.1400}}
  >{\raggedright\arraybackslash}p{(\linewidth - 10\tabcolsep) * \real{0.1400}}
  >{\raggedright\arraybackslash}p{(\linewidth - 10\tabcolsep) * \real{0.2400}}@{}}
\toprule\noalign{}
\begin{minipage}[b]{\linewidth}\raggedright
Symbol
\end{minipage} & \begin{minipage}[b]{\linewidth}\raggedright
Meaning
\end{minipage} & \begin{minipage}[b]{\linewidth}\raggedright
Thm 1
\end{minipage} & \begin{minipage}[b]{\linewidth}\raggedright
Thm 2
\end{minipage} & \begin{minipage}[b]{\linewidth}\raggedright
Thm 3
\end{minipage} & \begin{minipage}[b]{\linewidth}\raggedright
Definition
\end{minipage} \\
\midrule\noalign{}
\endhead
\bottomrule\noalign{}
\endlastfoot
\(\theta\) & Detection threshold & x & --- & --- &
\(\theta \in (0, 1)\) \\
\(NS(x)\) & Noise score & --- & x & --- & SCX noise score for sample
\(x\) \\
\(\hat{z}(x)\) & Predicted noise label & --- & x & --- &
\(\hat{z}(x) = \mathbf{1}\{NS(x) > t\}\) \\
\(\text{FPR}_s\) & False positive rate (state \(s\)) & x & --- & --- &
\(\mathbb{P}(C > \theta \mid \text{clean}, s)\) \\
\(\text{TPR}_s\) & True positive rate (state \(s\)) & x & --- & --- &
\(\mathbb{P}(C > \theta \mid \text{noise}, s)\) \\
F1 & F1 score of detector & x & x & --- &
\(2\text{TP}/(2\text{TP}+\text{FP}+\text{FN})\) \\
AUC & Area under ROC curve & --- & x & --- &
\(\mathbb{P}(\text{score}_{\text{noise}} > \text{score}_{\text{clean}})\) \\
PR-AUC & Area under precision-recall curve & --- & x & --- &
\(\int \text{Precision}(t) \, d\text{Recall}(t)\) \\
\end{longtable}

\subsubsection{1.5 Theorem-Specific
Parameters}\label{theorem-specific-parameters}

\begin{longtable}[]{@{}
  >{\raggedright\arraybackslash}p{(\linewidth - 6\tabcolsep) * \real{0.2000}}
  >{\raggedright\arraybackslash}p{(\linewidth - 6\tabcolsep) * \real{0.2250}}
  >{\raggedright\arraybackslash}p{(\linewidth - 6\tabcolsep) * \real{0.2750}}
  >{\raggedright\arraybackslash}p{(\linewidth - 6\tabcolsep) * \real{0.3000}}@{}}
\toprule\noalign{}
\begin{minipage}[b]{\linewidth}\raggedright
Symbol
\end{minipage} & \begin{minipage}[b]{\linewidth}\raggedright
Meaning
\end{minipage} & \begin{minipage}[b]{\linewidth}\raggedright
Appears In
\end{minipage} & \begin{minipage}[b]{\linewidth}\raggedright
Defined As
\end{minipage} \\
\midrule\noalign{}
\endhead
\bottomrule\noalign{}
\endlastfoot
\(\mathcal{P}_{\text{noise}}, \mathcal{P}_{\text{hard}}\) & Two
data-generating processes & Thm 3 & Constructive counterexample \\
\(\varepsilon_1, \varepsilon_2\) & Expert error rates in states
\(s_1, s_2\) & Thm 3 & \(\varepsilon_1 \in (0, 1/2)\),
\(\varepsilon_2 \in (\varepsilon_1, 1/2]\) \\
\(\tilde{P}\) & Auxiliary distribution (\(\phi \perp S\)) & Thm 2 &
\(\tilde{P}(\phi, S) = P(\phi)P(S)\) \\
\(F1_{\text{base}}\) & Loss-baseline detector F1 & Thm 2 &
\(\max_t \text{F1}(\hat{z}_{\text{loss}}(t))\) \\
\(F1_{\text{rand}}\) & Random detector F1 & Thm 2 &
\(2\eta/(1+\eta)\) \\
\end{longtable}

\begin{center}\rule{0.5\linewidth}{0.5pt}\end{center}

\subsection{2 Assumptions
Cross-Reference}\label{assumptions-cross-reference}

\begin{longtable}[]{@{}
  >{\raggedright\arraybackslash}p{(\linewidth - 8\tabcolsep) * \real{0.1111}}
  >{\raggedright\arraybackslash}p{(\linewidth - 8\tabcolsep) * \real{0.3056}}
  >{\raggedright\arraybackslash}p{(\linewidth - 8\tabcolsep) * \real{0.1944}}
  >{\raggedright\arraybackslash}p{(\linewidth - 8\tabcolsep) * \real{0.1944}}
  >{\raggedright\arraybackslash}p{(\linewidth - 8\tabcolsep) * \real{0.1944}}@{}}
\toprule\noalign{}
\begin{minipage}[b]{\linewidth}\raggedright
ID
\end{minipage} & \begin{minipage}[b]{\linewidth}\raggedright
Assumption
\end{minipage} & \begin{minipage}[b]{\linewidth}\raggedright
Thm 1
\end{minipage} & \begin{minipage}[b]{\linewidth}\raggedright
Thm 2
\end{minipage} & \begin{minipage}[b]{\linewidth}\raggedright
Thm 3
\end{minipage} \\
\midrule\noalign{}
\endhead
\bottomrule\noalign{}
\endlastfoot
A1 & Disjoint training sets & \textbf{Required} & Implicit & Target of
Corollary 2 \\
A2 & Clean-data conditional independence & \textbf{Required} & Implicit
& Technical basis for Cor. 2 \\
A3 & Bounded loss & \textbf{Required} & Implicit & --- \\
A4 & Uniform independent noise & \textbf{Required} & --- & Target of
Corollary 3 \\
A5 & State homogeneity & \textbf{Required} & --- & Target of Corollary
4 \\
A6 & Balanced error distribution & \textbf{Required} & --- & Target of
Corollary 5 \\
\end{longtable}

\begin{center}\rule{0.5\linewidth}{0.5pt}\end{center}

\subsection{3 Theorem Dependency Graph}\label{theorem-dependency-graph}

\begin{verbatim}
                         ┌──────────────────────┐
                         │   Theorem 3           │
                         │  (Unidentifiability)  │
                         │                       │
                         │  Core claim:          │
                         │  Without A1-A6,       │
                         │  noise vs. difficulty │
                         │  is observationally    │
                         │  equivalent           │
                         └──────────┬───────────┘
                                    │
                    "What minimum assumptions
                     are needed to make the
                     problem tractable?"
                                    │
                                    v
               ┌─────────────────────────────────────┐
               │   A1-A6 are the minimal sufficient   │
               │   conditions to break Thm 3's        │
               │   unidentifiability                   │
               └─────────────────────────────────────┘
                                    │
                    ┌───────────────┴───────────────┐
                    │                               │
                    v                               v
   ┌────────────────────────────┐    ┌────────────────────────────┐
   │   Theorem 1                │    │   Theorem 2                │
   │  (Noise Detection)         │    │  (Weak Feature Failure)    │
   │                            │    │                            │
   │  Positve result:           │    │  Negative result:          │
   │  Under A1-A6, noise is     │    │  Even under A1-A6, if      │
   │  detectable with F1→1      │    │  features are δ-weak,      │
   │  exponentially fast in M   │    │  SCX cannot beat loss      │
   │                            │    │  baseline by more than     │
   │                            │    │  O(√δ)                     │
   └────────────────────────────┘    └────────────────────────────┘
                    │                               │
                    └───────────────┬───────────────┘
                                    │
                                    v
                    ┌──────────────────────────────┐
                    │   Unified SCX Picture         │
                    │                               │
                    │  SCX Effectiveness ≈          │
                    │    I(φ; S) - O(1/√n)         │
                    │                               │
                    │  • Thm 3: Why we need A1-A6   │
                    │  • Thm 1: What we gain with   │
                    │    them                       │
                    │  • Thm 2: Where the boundary  │
                    │    lies                      │
                    └──────────────────────────────┘
\end{verbatim}

\subsubsection{3.1 Dependency Logic}\label{dependency-logic}

The three theorems form a logical chain of necessity, sufficiency, and
boundary:

\begin{enumerate}
\def\labelenumi{\arabic{enumi}.}
\item
  \textbf{Theorem 3 (Fundamental Hardness)} establishes that without
  structural assumptions, the noise detection problem is ill-posed:
  noise and intrinsic difficulty are observationally equivalent. This is
  the \textbf{necessity} result -- it explains \emph{why} assumptions
  are required, not just convenient.
\item
  \textbf{Theorem 1 (Sufficient Conditions)} shows that under
  assumptions A1-A6 -- the minimal set identified in Theorem 3 -- noise
  detection succeeds with exponentially decaying error. This is the
  \textbf{sufficiency} result -- it quantifies \emph{how well} SCX works
  when its assumptions hold.
\item
  \textbf{Theorem 2 (Boundary Condition)} shows that even when A1-A6
  hold, the feature representation \(\phi\) must carry enough mutual
  information about the true state \(S\) for SCX to outperform a simple
  loss baseline. This is the \textbf{boundary} result -- it identifies
  \emph{when} SCX ceases to add value despite valid assumptions.
\end{enumerate}

\subsubsection{3.2 Information-Theoretic
Unification}\label{information-theoretic-unification}

The three theorems can be unified under a single information-theoretic
inequality:

\[\underbrace{\text{SCX Advantage}}_{\text{Thm 1: exponential F1} \to 1} \;\leq\; \underbrace{f(I(\phi; S))}_{\text{Thm 2: } O(\sqrt{\delta}) \text{ bound}} \;\leq\; \underbrace{g(\text{A1-A6})}_{\text{Thm 3: minimal conditions}}\]

where \(f\) and \(g\) are monotone increasing functions. The chain
reads: - Without A1-A6, \(g = 0\) and the problem is unidentifiable (Thm
3) - With A1-A6 but \(I(\phi; S) \to 0\), SCX collapses to the baseline
(Thm 2) - With both A1-A6 and strong features, SCX achieves exponential
guarantee (Thm 1)

\begin{center}\rule{0.5\linewidth}{0.5pt}\end{center}

\subsection{4 Key Inequalities
Reference}\label{key-inequalities-reference}

\begin{longtable}[]{@{}
  >{\raggedright\arraybackslash}p{(\linewidth - 4\tabcolsep) * \real{0.3548}}
  >{\raggedright\arraybackslash}p{(\linewidth - 4\tabcolsep) * \real{0.2903}}
  >{\raggedright\arraybackslash}p{(\linewidth - 4\tabcolsep) * \real{0.3548}}@{}}
\toprule\noalign{}
\begin{minipage}[b]{\linewidth}\raggedright
Inequality
\end{minipage} & \begin{minipage}[b]{\linewidth}\raggedright
Used In
\end{minipage} & \begin{minipage}[b]{\linewidth}\raggedright
Statement
\end{minipage} \\
\midrule\noalign{}
\endhead
\bottomrule\noalign{}
\endlastfoot
Hoeffding & Thm 1 (Lemmas 2, 3) &
\(\mathbb{P}(\bar{X} - \mu \geq t) \leq e^{-2nt^2}\) \\
Chernoff (Sanov) & Thm 1 (tightened) &
\(\mathbb{P}(\bar{X} \geq \theta) \leq e^{-n \cdot \text{KL}(\theta \|\mu)}\) \\
Pinsker & Thm 2 (original) &
\(\operatorname{TV}(P, Q) \leq \sqrt{\frac{1}{2}\operatorname{KL}(P\|Q)}\) \\
Bretagnolle-Huber & Thm 2 (tightened) &
\(\operatorname{TV}(P, Q) \leq \sqrt{1 - e^{-\operatorname{KL}(P\|Q)}}\) \\
Fano & Thm 2 (Lemma 1) &
\(P(\hat{S} \neq S) \geq \frac{H(S \mid \phi) - \log 2}{\log \|\mathcal{S}\|}\) \\
Data Processing & Thm 2 &
\(\operatorname{TV}(P_{\text{pred}}, Q_{\text{pred}}) \leq \operatorname{TV}(P, Q)\) \\
\end{longtable}

\begin{center}\rule{0.5\linewidth}{0.5pt}\end{center}

\subsection{5 Quick-Start Guide}\label{quick-start-guide}

\begin{longtable}[]{@{}
  >{\raggedright\arraybackslash}p{(\linewidth - 2\tabcolsep) * \real{0.5686}}
  >{\raggedright\arraybackslash}p{(\linewidth - 2\tabcolsep) * \real{0.4314}}@{}}
\toprule\noalign{}
\begin{minipage}[b]{\linewidth}\raggedright
If you want to understand\ldots{}
\end{minipage} & \begin{minipage}[b]{\linewidth}\raggedright
Read this combination
\end{minipage} \\
\midrule\noalign{}
\endhead
\bottomrule\noalign{}
\endlastfoot
Why SCX needs assumptions & Thm 3 + Section 4 of Thm 3 (A1-A6
mapping) \\
How well SCX detects noise & Thm 1 + Corollaries 1-4 \\
When SCX cannot work & Thm 2 + Diagnostic Section 7 \\
Practical deployment & Thm 1 Corollary 4 (finite-sample) + Thm 2
diagnostics \\
The complete picture & All three + this notation document \\
\end{longtable}

\begin{center}\rule{0.5\linewidth}{0.5pt}\end{center}

\textbf{Document version}: 2026-06-27

\end{document}
